\chapter*{数列}
\subsection*{等差中项、等比中项}
\begin{enumerate}
    \item (2025$\cdot$浦东新区一模$\cdot$7)若等差数列 $\left\{  {a}_{n}\right\}$ 满足 ${a}_{7} + {a}_{8} + {a}_{9} = 0,{a}_{7} + {a}_{10} = 1$ ，则 ${a}_{1} =$ \blkda[2]{$-7$}.
    \begin{jieda}
        因为数列 $\left\{  {a}_{n}\right\}$ 为等差数列,则 ${a}_{7} + {a}_{8} + {a}_{9} = 3{a}_{8} =$ 0,即 ${a}_{8} = 0$ ,且 ${a}_{7} + {a}_{10} = {a}_{8} + {a}_{9} = 1$ ,可得 ${a}_{9} = 1$ ,即 $\left\{  \begin{array}{l} {a}_{8} = {a}_{1} + {7d} = 0 \\  {a}_{9} = {a}_{1} + {8d} = 1 \end{array}\right.$ ,解得 $\left\{  \begin{array}{l} {a}_{1} =  - 7 \\  d = 1 \end{array}\right.$
    \end{jieda}
    \item (2025$\cdot$黄浦区二模$\cdot$11)设 $\left\{  {a}_{n}\right\}$ 为等差数列，其前 $n$ 项和为 ${S}_{n}$ ，若 $\left( {{S}_{8} - {S}_{7}}\right)({S}_{9} - {S}_{7}) < 0$ ，则满足 ${S}_{m}{S}_{m + 1} < 0$ 的正整数 $m =$ \blkda[2]{$15$}.
    \begin{jieda}
        $a_8 \cdot \left(a_8+a_9\right) < 0$，不妨设公差$d < 0$，则$a_8 > 0, a_8+a_9 < 0$，故$S_{15} > 0, S_{16} < 0$，因此$m = 15$.
    \end{jieda}
    \item (2025$\cdot$松江区一模$\cdot$8)已知数列 $\left\{  {a}_{n}\right\}$ 是等比数列，若 ${\log }_{2}{a}_{1} + {\log }_{2}{a}_{4} = 3$ ， ${{2}}^{{a}_{2}} \cdot  {2}^{{a}_{3}} = {64}$ ，则 ${a}_{10} =$ \blkda[2]{512 或 $\dfrac{1}{64}$}.
    \begin{jieda}
        因为 ${\log }_{2}{a}_{1} + {\log }_{2}{a}_{4} = {\log }_{2}\left( {{a}_{1} \cdot  {a}_{4}}\right)  = 3$ ,所以 ${a}_{1} \cdot  {a}_{4} = {a}_{2} \cdot  {a}_{3} = 8$ ,因为 ${2}^{{a}_{2}} \cdot  {2}^{{a}_{3}} = {2}^{{a}_{2} + {a}_{3}} = {64}$ ,所以 ${a}_{2} + {a}_{3} =$ 6,解得 ${a}_{2} = 2,{a}_{3} = 4$ 或 ${a}_{2} = 4,{a}_{3} = 2$ .

        \ding{172}  当 ${a}_{2} = 2,{a}_{3} = 4$ 时， $q = 2 \Longrightarrow  {a}_{10} = 2 \times  {2}^{8} = {512}$ ；

        \ding{173}  当 ${a}_{2} = 4,{a}_{3} = 2$ 时， $q = \dfrac{1}{2} \Longrightarrow  {a}_{10} = 4 \times  {\left( \dfrac{1}{2}\right) }^{8} = \dfrac{1}{64}$ .
    \end{jieda}
    \item 已知$a = 243, b = 3$，$c$是$a, b$的等比中项
    \begin{enumerate}
        \item 若$d, e$分别是$a, c$和$c, b$的等差中项，则$a+b+c+d+e = $\blkda[4]{423或315}
        \item 若$d, e$分别是$a, c$和$c, b$的等比中项，则$a+b+c+d+e = $\blkda[4]{363或345或201或183}
    \end{enumerate}
    \begin{jieda}
        首先可以得到$c = -27$或27.
        \begin{enumerate}
            \item $a+b+c+d+e = a+b+c + \dfrac{a+c}{2} + \dfrac{b+c}{2} = \dfrac{3(a+b)}{2} + 2c = $ 423或315
            \item 此时必有$c = 27$，\begin{enumerate}[1.]
                \item $d = 81, e = 9$，$a+b+c+d+e = 363$
                \item $d = 81, e = -9$，$a+b+c+d+e = 345$
                \item $d = -81, e = 9$，$a+b+c+d+e = 201$
                \item $d = -81, e = -9$，$a+b+c+d+e = 183$
            \end{enumerate}
        \end{enumerate}
    \end{jieda}
\end{enumerate}

\subsection*{等差中项：$(2n-1)a_n = S_{2n-1}$}
\begin{enumerate}
    \item 记两个等差数列$\{a_n\}, \{b_n\}$的前$n$项和分别是$S_n$和$T_n$，且$\dfrac{S_n}{T_n} = \dfrac{7n+1}{4n+27}(n \in \mathbf{N^*})$，则$\dfrac{a_{11}}{b_{11}} = $\blkda{$\dfrac{4}{3}$}
    \begin{jieda}
        $\dfrac{a_{11}}{b_{11}} = \dfrac{21 \cdot a_{11}}{21 \cdot b_{11}} = \dfrac{S_{21}}{T_{21}} = \dfrac{148}{111} = \dfrac{4}{3}$.

        关键结论：$(2n-1)a_n = S_{2n-1}$
    \end{jieda}
    \item 已知等差数列 $\left\{  {a}_{n}\right\}$ 和等差数列 $\left\{  {b}_{n}\right\}$ 的前 $n$ 项和分别为 ${S}_{n}$ 和 ${T}_{n}$ ,且 $\dfrac{{S}_{n}}{T_n} = \dfrac{{5n} + {63}}{n + 3}$ ,则使得 $\dfrac{a_n}{b_n}$整数的正整数 $n$ 的个数为 \dotfill{(\kh{B})}
    \xx{6}{7}{8}{9}
    \begin{jieda}
        由于 ${S}_{{2n} - 1} = \dfrac{\left( {{a}_{1} + {a}_{{2n} - 1}}\right) \left( {{2n} - 1}\right) }{2} = \dfrac{2{a}_{n}\left( {{2n} - 1}\right) }{2} = \left( {{2n} - 1}\right) {a}_{n}$
        
        所以 $\dfrac{{a}_{n}}{{b}_{n}} = \dfrac{{S}_{{2n} - 1}}{{T}_{{2n} - 1}} = \dfrac{5\left( {{2n} - 1}\right)  + {63}}{{2n} - 1 + 3} = \dfrac{{5n} + {29}}{n + 1} = 5 + \dfrac{24}{n + 1}$ 
        
        要使 $\dfrac{{a}_n}{b_n}$ 为整数,则 $n + 1$ 为 24 的因数,由于 $n + 1 \geq  2$ ,故 $n + 1$ 可以为2,3,4,6,8,12,24,故满足条件的正整数 $n$ 的个数为 7 个
    \end{jieda}
\end{enumerate}

\subsection*{等差数列、等比数列的通项公式}
\begin{enumerate}
    \item (2025$\cdot$金山区一模$\cdot$8)已知 ${S}_{n}$ 是等差数列 $\left\{  {a}_{n}\right\}$ 的前 $n$ 项和，若 $2{a}_{1} + 3{a}_{11} = {20}$ ，则 ${S}_{13}$ 的值为\blkda[2]{52}.
    \begin{jieda}
        设等差数列 $\left\{  {a}_{n}\right\}$ 的公差为 $d$ ,则由题得 $2{a}_{1} + 3{a}_{11} =  5{a}_{1} + {30d} = {20}$ 即 ${a}_{1} + {6d} = 4$ ,所以 ${S}_{13} = {13}{a}_{1} + \dfrac{{13} \times  {12}}{2}d =  {13}\left( {{a}_{1} + {6d}}\right)  = {13} \times  4 = {52}$.
    \end{jieda}
    \item (2025$\cdot$长宁区一模$\cdot$16)数列 $\left\{  {a}_{n}\right\}$ 为严格增数列，且对任意的正整数 $n$ ，都有 $\dfrac{{a}_{n + 1}}{n + 1} \geq  \dfrac{{a}_{n}}{n}$ ，则称数列 $\left\{  {a}_{n}\right\}$ 满足“性质 $\Omega$ ”. 现给出两个命题:
    \ding{172}  存在等差数列 $\left\{  {a}_{n}\right\}$ 满足“性质 $\Omega$ ”；

    \ding{173}  任意等比数列 $\left\{  {a}_{n}\right\}$ ，若首项 ${a}_{1} > 0$ ，则 $\left\{  {a}_{n}\right\}$ 满足“性质 $\Omega$ ”.

    下列判断正确的是 \dotfill{(\kh{B})}

    \xx{\ding{172}  \ding{173} 都是真命题}{\ding{172} 是真命题， \ding{173} 是假命题}{\ding{172} 是假命题， \ding{173} 是真命题}{\ding{172}  \ding{173} 都是假命题}
    \begin{jieda}
        设等差数列 $\left\{  {a}_{n}\right\}$通项公式为$a_n = an+b, a > 0$，则满足“性质 $\Omega$ ”即$\dfrac{a(n+1)+b}{n+1} \geq \dfrac{an+b}{n}$，当$b = 0$时即可.

        设等比数列 $\left\{  {a}_{n}\right\}$通项公式为$a_n = a_1q^{n-1}, a_1 > 0, q > 1$，则满足“性质 $\Omega$ ”即$\dfrac{a_1q^n}{n+1} \geq \dfrac{a_1q^{n-1}}{n}$，即$q \geq \dfrac{n+1}{n}$，故$q \geq 2$.

        \ps 举反例$q = 1$也可.
    \end{jieda}
    \item (2024$\cdot$上海秋季卷$\cdot$12)无穷等比数列 $\left\{  {a}_{n}\right\}$ 满足首项 ${a}_{1} > 0,q > 1$ ，记 ${I}_{n} = \left\{  {x - y \mid  x,y \in  \left\lbrack  {{a}_{1},{a}_{2}}\right\rbrack   \cup  }\right.  \left. \left\lbrack  {{a}_{n},{a}_{n + 1}}\right\rbrack  \right\}$ ,若对任意正整数 $n$ 集合 ${I}_{n}$ 是闭区间,则 $q$ 的取值范围是\blkda[2]{$\lbrack 2, + \infty)$}.
    \begin{jieda}
        根据${a}_{1} > 0,q > 1$可知$a_n$严增，不妨设$x > y$

        故$I_n = [0, a_2-a_1] \cup [0, a_{n+1} - a_n] \cup [a_n-a_2, a_{n+1}-a_1]$

        易知$a_{n+1} - a_n \geq a_2-a_1$，由$I_n$为闭区间可知$a_{n+1} - a_n \geq a_n-a_2$，即$a_1q^n+a_1q \geq 2a_1q^{n-1}$，故$q^n+q \geq 2q^{n-1}$，即$q + \dfrac{1}{q^{n-2}} \geq 2$恒成立，左侧随着$n$增加严减，故考虑$n \rightarrow +\infty$时，故$q \geq 2$
    \end{jieda}
\end{enumerate}
\subsection*{等差数列、等比数列$n$项和}
\begin{enumerate}
    \item (2024$\cdot$松江区二模$\cdot$7)已知等差数列 $\left\{  {a}_{n}\right\}$ 的公差为 2，前 $n$ 项和为 ${S}_{n}$ ，若 ${a}_{3} = {S}_{5}$ ，则使得 ${S}_{n} < {a}_{n}$ 成立的 $n$ 的最大值为\blkda[2]{5}.
    \begin{jieda}
        ${a}_{1} + {2d} = 5{a}_{1} + \dfrac{5 \times  4}{2} \times  d \Longrightarrow  {a}_{1} =  - 4,{S}_{n} = {n}^{2} - {5n}$ , ${a}_{n} = {2n} - 6$ ,由 ${S}_{n} < {a}_{n}$ ,得 ${n}^{2} - {7n} + 6 < 0$ ,解得 $1 < n < 6$ ,故 $n$ 的最大值为 5.
    \end{jieda}
    \item (2024$\cdot$黄浦区二模$\cdot$10)已知数列 $\left\{  {a}_{n}\right\}$ 是给定的等差数列，其前 $n$ 项和为 ${S}_{n}$ ，若 ${a}_{9}{a}_{10} < 0$ ，且当 $m =  {m}_{0}$ 与 $n = {n}_{0}$ 时， $\left| {{S}_{m} - {S}_{n}}\right|$ ( $m,n \in  \left\{  {x \mid  x \leq  {30},x \in  {\mathbf{N}}^{ * }}\right\}$ )取得最大值，则 $\left| {{m}_{0} - {n}_{0}}\right|$ 的值为\blkda[2]{21}.
    \begin{jieda}
        设等差数列 $\left\{  {a}_{n}\right\}$ 的公差为 $d$ . 当 $d > 0$ 时,由于 ${a}_{9}{a}_{10} <$ 0,所以 ${a}_{9} < 0 < {a}_{10}$ ,数列 $\left\{  {a}_{n}\right\}$ 为严格增数列,故 ${S}_{n}$ 在区间 $\lbrack 1$ , 30]上， ${S}_{9}$ 为最小值， ${S}_{30}$ 为最大值，若 $\left| {{S}_{m} - {S}_{n}}\right|$ 取得最大值， 则 $m = {30},n = 9$ ,故 $\left| {{m}_{0} - {n}_{0}}\right|  = \left| {{30} - 9}\right|  = {21}$ ; 当 $d < 0$ 时, ${a}_{9} > 0 > {a}_{10}$ ,数列 $\left\{  {a}_{n}\right\}$ 为严格减数列, ${S}_{n}$ 在 $n \in  \left\lbrack  {1,{30}}\right\rbrack  ,{S}_{9}$ 为最大值, ${S}_{30}$ 为最小值,若 $\left| {{S}_{m} - {S}_{n}}\right|$ 取得最大值,则 $m = 9,n =$ 30，故 $\left| {{m}_{0} - {n}_{0}}\right|  = \left| {9 - {30}}\right|  = {21}$ ，综上， $\left| {{m}_{0} - {n}_{0}}\right|  = {21}$ .
    \end{jieda}
    \item (2025$\cdot$黄浦区二模$\cdot$7)已知等比数列 $\left\{  {a}_{n}\right\}$ 为严格增数列，其前 $n$ 项和为 ${S}_{n}$ ，若 ${a}_{1}{a}_{10} = {a}_{8}$ ， ${S}_{4} - {S}_{1} = \dfrac{21}{4}$ ，则该数列的公比为\blkda[2]{4}.
    \begin{jieda}
        设等比数列 $\left\{  {a}_{n}\right\}$ 的公比为 $q$ ,由 ${a}_{1}{a}_{10} = {a}_{8}$ ,则 ${a}_{1}{a}_{1}{q}^{9} = {a}_{1}{q}^{7}$ ,解得 ${a}_{1} = {q}^{-2}$ ,由 ${S}_{4} - {S}_{1} = \dfrac{21}{4}$ ,则 $\dfrac{{a}_{1}\left( {1 - {q}^{4}}\right) }{1 - q} -  {a}_{1} = \dfrac{21}{4}$ ,将 ${a}_{1} = {q}^{-2}$ 代入上式可得 $\dfrac{1 - {q}^{4}}{\left( {1 - q}\right) {q}^{2}} - \dfrac{1}{{q}^{2}} = \dfrac{21}{4}$ ,去分母可得 $4{q}^{4} - {21}{q}^{3} + {21}{q}^{2} - {4q} = 0$ ,易知 $q \neq  0$ ,可得 $4{q}^{3} - {21}{q}^{2}  + {21q} - 4 = 0$ ,分解因式可得 $\left( {q - 1}\right) \left( {{4q} - 1}\right) \left( {q - 4}\right)  = 0$ ,易知 $q \neq  1$ ,解得 $q = \dfrac{1}{4}$ 或 4,当 $q = \dfrac{1}{4}$ 时, ${a}_{1} = {16}$ ,则 ${a}_{n} = {16} \times  {\left( \dfrac{1}{4}\right) }^{n - 1} = {4}^{3 - n}$ ,单调递减,不合题意.
    \end{jieda}
    \item (2025$\cdot$青浦区二模$\cdot$7)数列 $\left\{  {a}_{n}\right\}$ 中， ${a}_{1} = 1$ ， ${a}_{n} = 2{a}_{n + 1}$ ，则 $\mathop{\sum }\limits_{{i = 1}}^{{+\infty }}{a}_{i} =$ \blkda[2]{2}.
    \begin{jieda}
        ${a}_{1} = 1,{a}_{n} = 2{a}_{n + 1}$ ,即 $\dfrac{{a}_{n + 1}}{{a}_{n}} = \dfrac{1}{2},\left\{  {a}_{n}\right\}$ 是等比数列, 公比为 $\dfrac{1}{2}$ ,所以 $\mathop{\sum }\limits_{{i = 1}}^{{+\infty }}{a}_{i} = \dfrac{1}{1 - \dfrac{1}{2}} = 2$.
    \end{jieda}
    \item (2024$\cdot$闵行区一模$\cdot$11)已知数列 $\left\{  {a}_{n}\right\}$ 为无穷等比数列，若 $\mathop{\sum }\limits_{{i = 1}}^{{+\infty }}{a}_{i} =  - 2$ ，则 $\mathop{\sum }\limits_{{i = 1}}^{{+\infty }}\left| {a}_{i}\right|$ 的取值范围为\blkda[2]{$[2, + \infty)$}.
    \begin{jieda}
        因为 $\left\{  {a}_{n}\right\}$ 为无穷等比数列, $\mathop{\sum }\limits_{{i = 1}}^{{+\infty }}{a}_{i} =  - 2$ ,所以 $0 <  \left| q\right|  < 1$ ,则 $\mathop{\sum }\limits_{{i = 1}}^{{+\infty }}{a}_{i} = \dfrac{{a}_{1}}{1 - q} =  - 2$ ,则 ${a}_{1} =  - 2\left( {1 - q}\right)$ ,因为 $\dfrac{\left| {a}_{n + 1}\right| }{\left| {a}_{n}\right| } = \left| \dfrac{{a}_{n + 1}}{{a}_{n}}\right|  = \left| q\right|$ ,所以 $\left| {a}_{n}\right|$ 是以 $\left| q\right|$ 为公比的等比数列,且 $0 < \left| q\right|  < 1$ ,此时 $1 - |q| > 0$ ,所以 $\mathop{\sum }\limits_{{i = 1}}^{{+\infty }}\left| {a}_{i}\right|  = \dfrac{\left| {a}_{1}\right| }{1 - \left| q\right| } =  \dfrac{\left| -2\left( 1 - q\right) \right| }{1 - \left| q\right| } = \dfrac{2\left( {1 - q}\right) }{1 - \left| q\right| }$
        
        当 $0 < q < 1$ 时, $\dfrac{2\left( {1 - q}\right) }{1 - \left| q\right| } = \dfrac{2\left( {1 - q}\right) }{1 - q}  = 2$
        
        当 $- 1 < q < 0$ 时, $\dfrac{2\left( {1 - q}\right) }{1 - \left| q\right| } = \dfrac{2\left( {1 - q}\right) }{1 + q} = \dfrac{4 - 2\left( {1 + q}\right) }{1 + q} =  \dfrac{4}{1 + q} - 2$ ,因为 $- 1 < q < 0$ ,所以 $0 < 1 + q < 1$ ,故 $\dfrac{4}{1 + q} > 4$ ,则 $\dfrac{4}{1 + q} - 2 > 2$
        
        综上: $\dfrac{2\left( {1 - q}\right) }{1 - \left| q\right| } \geq  2$ ,即 $\mathop{\sum }\limits_{{i = 1}}^{{+\infty }}\left| {a}_{i}\right|  \geq  2$ ,故 $\mathop{\sum }\limits_{{i = 1}}^{{+\infty }}\left| {a}_{i}\right|$ 的取值范围为 $\lbrack 2, + \infty )$ .
    \end{jieda}
\end{enumerate}
\subsection*{等差数列、等比数列前$n$项和的形式}
\begin{enumerate}
    \item (2024$\cdot$静安区二模$\cdot$7)已知等比数列的前 $n$ 项和为 ${S}_{n} = {\left( \dfrac{1}{2}\right) }^{n} + a$ ，则 $a$ 的值为\blkda[2]{$-1$}.
    \begin{jieda}
        由等比数列的前 $n$ 项和为 ${S}_{n} = A - A \cdot q^n$，故$\left\{\begin{array}{l}
            A = a \\
            -A \cdot q^n = {\left( \dfrac{1}{2}\right) }^{n}
        \end{array}\right.$，故$\left\{\begin{array}{l}
            q = \dfrac{1}{2} \\
            A = -1.
        \end{array}\right.$
    \end{jieda}
    \item 记 ${S}_{n}$ 为数列 $\left\{  {a}_{n}\right\}$ 的前 $n$ 项和,设甲: $\left\{  {a}_{n}\right\}$ 为等差数列; 乙: $\left\{  \dfrac{{S}_{n}}{n}\right\}$ 为等差数列, 则甲是乙的{(\kh{C})}
    \xx{充分非必要条件}{必要非充分条件}{充要条件}{既非充分又非必要条件}
    \item (2026$\cdot$奉贤区一模$\cdot$16)已知等差数列 $\left\{  {a}_{n}\right\}$ 的公差不为零, ${S}_{n}$ 为其前 $n$ 项和,存在正整数 $k$ 满足 ${S}_{k} = 0$ ,有两个命题:
    
    命题\ding{172}: 设数列公差 $d$ ,则 $d \cdot  {S}_{1} < 0$ .
    
    命题\ding{173}: $i, j$ 均是小于 $k$ 的正整数,则 ${S}_{i} \cdot  {S}_{j} = {S}_{k - i}{S}_{k - j}$ .
    
    以上判断正确的是\dotfill{(\kh{D})}
    \xx{\ding{172}  \ding{173} 均是真命题}{\ding{172} 是真命题， \ding{173} 是假命题}{\ding{172}  \ding{173} 均是假命题}{\ding{172} 是假命题， \ding{173} 是真命题}
    \begin{jieda}
        \begin{dingautolist}{172}
            \item 若$k = 1$，则$d$的正负判断不了.
            \item 根据$S_n$是过原点的二次可知$S_i = S_{k-i}$，$S_j = S_{k-j}$
        \end{dingautolist}
    \end{jieda}
    \item (2025$\cdot$徐汇区一模$\cdot$16)已知数列 $\left\{  {a}_{n}\right\}$ 的前 $n$ 项和为 ${S}_{n}$ ，设 ${t}_{n} = \dfrac{{S}_{n}}{n}$ ( $n$ 为正整数). 若存在常数 $c$ ， 使得任意两两不相等的正整数 $i,j,k$ ,都有 $\left( {i - j}\right) {t}_{k} + \left( {j - k}\right) {t}_{i} + \left( {k - i}\right) {t}_{j} = c$ ,则称数列 $\left\{  {a}_{n}\right\}$ 为 “轮换均值数列”. 现有下列两个命题: \ding{172}  任意等差数列 $\left\{  {a}_{n}\right\}$ 都是“轮换均值数列”， \ding{173}  存在公比不为 1 的等比数列 $\left\{  {b}_{n}\right\}$ 是“轮换均值数列”. 则下列说法正确的是 \dotfill{(\kh{A})}
    \xx{\ding{172} 是真命题， \ding{173} 是假命题}{\ding{172} 是假命题， \ding{173} 是真命题}{\ding{172}  \ding{173} 都是真命题}{\ding{172}  \ding{173} 都是假命题}
    \begin{jieda}
        对于等差数列 $\left\{  {a}_{n}\right\}$可设$t_n = pn+q$，显然满足要求.

        对于公比不为 1 的等比数列 $\left\{  {b}_{n}\right\}$可设$t_n = \dfrac{A}{n}(1-q^n)$，显然不满足要求.
    \end{jieda}

    \item (2024$\cdot$长宁区二模$\cdot$16)设数列 $\left\{  {a}_{n}\right\}$ 的前 $n$ 项和为 ${S}_{n}$ ，若存在非零常数 $c$ ，使得对任意正整数 $n$ ， 都有 $2\sqrt{{S}_{n}} = {a}_{n} + c$ ,则称数列 $\left\{  {a}_{n}\right\}$ 具有性质 $p$. 现有如下两个命题: \ding{172}  存在等差数列 $\left\{  {a}_{n}\right\}$ 具有性质 $p$ ;  \ding{173}  不存在等比数列 $\left\{  {a}_{n}\right\}$ 具有性质 $p$. 下列判断正确的是 \dotfill{(\kh{B})}
    \xx{\ding{172}  \ding{173} 均是真命题}{\ding{172} 是真命题， \ding{173} 是假命题}{\ding{172} 是假命题， \ding{173} 是真命题}{\ding{172}  \ding{173} 均是假命题}
    \begin{jieda}
        根据等差数列前$n$项和的形式可知$S_n = k\cdot n^2$可以满足要求，具体计算可知对 ${a}_{n} = {2n} - 1$ ,知 $\left\{  {a}_{n}\right\}$ 是等差数列,而 ${S}_{n} =  \dfrac{1}{2}n \cdot  \left( {1 + {2n} - 1}\right)  = {n}^{2}$ ,令 $c = 1$ 就有 $2\sqrt{{S}_{n}} = {2n} = {2n} -  1 + 1 = {a}_{n} + c$ ,所以 $\left\{  {a}_{n}\right\}$ 具有性质 $p$ ,这表明存在等差数列 $\left\{  {a}_{n}\right\}$ 具有性质 $p$
        
        根据等比数列前$n$项和的形式可知$q = -1$可以满足要求，故根据前两项有$\left\{\begin{array}{l}
            2\sqrt{a_1} = a_1+c \\
            0 = -a_1+c
        \end{array}\right.$，解得$\left\{\begin{array}{l}
            a_1 = 1 \\
            c = 1
        \end{array}\right.$，经检验$a_n = (-1)^{n-1}$满足要求，所以$\left\{  {a}_{n}\right\}$ 具有性质 $p$ ,这表明存在等比数列 $\left\{  {a}_{n}\right\}$ 具有性质 $p$ . 综上, \ding{172}  真、 \ding{173} 假. 故选: B.
    \end{jieda}
    \item (2024$\cdot$奉贤区一模$\cdot$16)已知等差数列 $\left\{  {a}_{n}\right\}$ 的前 $n$ 项和为 ${S}_{n}$ ，且关于正整数 $n$ 的不等式($S_n - {2022})\left( {{S}_{n + 1} - {2022}}\right)  < 0$ 与不等式 $\left( {{S}_{n} - {2023}}\right) \left( {{S}_{n + 1} - {2023}}\right)  < 0$ 的解集均为 $M$.

    命题 $\alpha$ : 集合 $M$ 中元素的个数一定是偶数个;

    命题 $\beta$ : 若数列 $\left\{  {a}_{n}\right\}$ 的公差 $d > 0$ ,且 ${n}_{0} \in  M$ ,则 ${a}_{{n}_{0} + 1} > 1$.

    下列说法中正确的是 \dotfill{(\kh{B})}

    \xx{命题 $\alpha$ 是真命题，命题 $\beta$ 是假命题}{命题 $\alpha$ 是假命题，命题 $\beta$ 是真命题}{命题 $\alpha$ 、命题 $\beta$ 都是假命题}{命题 $\alpha$ 、命题 $\beta$ 都是真命题} 
    \begin{jieda}
        命题 $\alpha  : {a}_{n} = 4,{S}_{n} = {4n},{S}_{n + 1} = {4n} + 4$ ,

        $\left\{  \begin{array}{l} \left( {{S}_{n} - {2022}}\right) \left( {{S}_{n + 1} - {2022}}\right)  < 0 \Longrightarrow  \left( {{4n} - {2022}}\right) \left( {{4n} - {2018}}\right)  < 0 \\  \left( {{S}_{n} - {2023}}\right) \left( {{S}_{n + 1} - {2023}}\right)  < 0 \Longrightarrow  \left( {{4n} - {2023}}\right) \left( {{4n} - {2019}}\right)  < 0 \end{array}\right.$，$n \in  \mathbf{N} \Longrightarrow  n = {505}$ ,解集均为 $M = \{ {505}\}$ ,唯一,故命题 $\alpha$ 为假命题; 

        \ps 根据等差数列前$n$项和的形式易知，一次和二次的形式均可以举出反例，上面是一次的反例，二次的情况下对称轴位置顶点纵坐标大于2023，两个第二大的纵坐标小于2022即可.
        
        命题 $\beta$ : 若数列 $\left\{  {a}_{n}\right\}$ 的公差 $d > 0$，则有$\left\{  {\begin{array}{l} {S}_{{n}_{0}} < {2022} \\  {S}_{{n}_{0} + 1} > {2023} \end{array} \Longrightarrow  {a}_{{n}_{0} + 1} = }\right. {S}_{{n}_{0} + 1} - {S}_{{n}_{0}} > 1$ ,命题 $\beta$ 是真命题; 综上所述: 命题 $\alpha$ 是假命题,命题 $\beta$ 是真命题.
    \end{jieda}
\end{enumerate}
\subsection*{$S_n$的不等关系}
\begin{enumerate}
    \item 设 $\left\{  {a}_{n}\right\}$ 是等差数列, ${S}_{n}$ 是其前 $n$ 项和,且 ${S}_{5} < {S}_{6},{S}_{6} = {S}_{7} > {S}_{8}$ ,则下列结论错误的是 \dotfill{(\kh{C})}
    \xx{$d < 0$}{${a}_{7} = 0$}{${S}_{9} > {S}_{5}$}{${S}_{6}$ 与 ${S}_{7}$ 均为 ${S}_{n}$ 的最大值}
    \begin{jieda}
        ${S}_{5} < {S}_{6},{S}_{6} = {S}_{7} > {S}_{8} \Longleftrightarrow a_6 > 0, a_7 = 0, a_8 < 0$，易知$A,B,D$正确

        $S_5 = S_8 > S_9$，C错误
    \end{jieda}
    \item 已知等差数列 $\left\{  {a}_{n}\right\}$ 的前 $n$ 项和为 ${S}_{n},$ 若 ${S}_{23} > 0,{S}_{24} < 0$ ,则下列结论正确的是 \dotfill{(\kh{D})}
    \xx{数列 $\left\{  {a}_{n}\right\}$ 是递增数列}{${a}_{13} > 0$}{当 ${S}_{n}$ 取得最大值时, $n = {13}$}{$\left| {a}_{13}\right|  > \left| {a}_{12}\right|$}
    \item (2024$\cdot$青浦区二模$\cdot$15)设 ${S}_{n}$ 是首项为 ${a}_{1}$ 、公比为 $q$ 的等比数列 $\left\{  {a}_{n}\right\}$ 的前 $n$ 项和，且 ${S}_{2023} <  {S}_{2025} < {S}_{2024}$ ,则 \dotfill{(\kh{C})}
    \xx{${a}_{1} > 0$}{$q > 0$}{$\left| {S}_{n}\right|  \leq  \left| {a}_{1}\right|$}{$\left| {S}_{n}\right|  < \left| q\right|$}
    \begin{jieda}
        因为 ${S}_{2023} < {S}_{2025} < {S}_{2024}$ ,所以 ${S}_{2025} - {S}_{2023} > 0$ , ${S}_{2025} - {S}_{2024} < 0$ ,即 ${a}_{2024} + {a}_{2025} > 0$ 且 ${a}_{2025} < 0$ ,所以 ${a}_{2024} >  - {a}_{2025} > 0$ ,且 ${a}_{2024} > 0$ ,两边都除以 ${a}_{2024}$ ,得 $1 >  - q > 0$ ,可得 $- 1 < q < 0$.
        
        对于 ${A}$ ,由 ${a}_{2025} = {a}_{1}{q}^{2024} < 0$ ,可得 ${a}_{1} < 0$ ,故 A 项不正确; 对于 B,由于 $- 1 < q < 0$ ,所以 $q > 0$ 不成立,故 B 不正确; 对于 ${C}$ ,因为 $- 1 < q < 0$ ,所以 $0 < 1 - {q}^{n} \leq  1 - q$ ,可得 $0 <  \dfrac{1 - {q}^{n}}{1 - q} \leq  1$ . 结合 ${S}_{n} = \dfrac{{a}_{1}\left( {1 - {q}^{n}}\right) }{1 - q}$ ,可得 $\left| {S}_{n}\right|  = \left| {a}_{1}\right|  \cdot  \left| \dfrac{1 - {q}^{n}}{1 - q}\right|  \leq  \left| {a}_{1}\right|$ ,故 C 正确; 对于 ${D}$ ,根据 $- 1 < q < 0$ 且 ${a}_{1} < 0$ ,当 ${a}_{1} =  - 1,q =  - \dfrac{1}{2}$ 时, $\left| {S}_{1}\right|  = 1 > \left| q\right|  = \dfrac{1}{2}$ ,此时 $\left| {S}_{n}\right|  < \left| q\right|$ 不成立，故 D 不正确. 故选:C.
    \end{jieda}
\end{enumerate}
\subsection*{$a_n = S_n - S_{n-1}$}
\begin{enumerate}
    \item (2025$\cdot$青浦区一模$\cdot$6)已知数列 $\left\{  {a}_{n}\right\}$ 满足 ${a}_{1} + 2{a}_{2} + 3{a}_{3} + \cdots  + n{a}_{n} = n\left( {n + 2}\right)$ ，则 ${a}_{66} =$ \blkda[2]{$\dfrac{133}{66}$}.
    \begin{jieda}
        因为 ${a}_{1} + 2{a}_{2} + 3{a}_{3} + \cdots  + n{a}_{n} = n\left( {n + 2}\right)$  \ding{172} ,当 $n \geq  2$ 时, ${a}_{1} + 2{a}_{2} + 3{a}_{3} + \cdots  + \left( {n - 1}\right) {a}_{n - 1} = \left( {n - 1}\right) \left( {n + 1}\right)$  \ding{173} ,  \ding{172} - \ding{173} 得 $n{a}_{n} = {2n} + 1$ ，所以 ${a}_{n} = \dfrac{{2n} + 1}{n}\left( {n \geq  2}\right)$ ，所以 ${a}_{66} =  \dfrac{133}{66}$.
    \end{jieda}
    \item (2025$\cdot$宝山区二模$\cdot$16)若对任意正整数 $n$ ，数列 $\left\{  {a}_{n}\right\}$ 的前 $n$ 项和 ${S}_{n}$ 都是完全平方数，则称数列 $\left\{  {a}_{n}\right\}$ 为“完全平方数列”. 有如下两个命题:  \ding{172}  若数列 $\left\{  {b}_{n}\right\}$ 的前 $n$ 项和 ${T}_{n} = {\left( n - t\right) }^{2}$ ( $t$ 为正整数),则使得数列 $\left\{  \left| {b}_{n}\right| \right\}$ 为“完全平方数列”的 $t$ 值有且仅有一个;  \ding{173}  存在无穷多个“完全平方数列”的等差数列. 对于以上两个命题, 下列选项中正确的是 \dotfill{(\kh{A})}

    \xx{\ding{172}  \ding{173} 都是真命题}{\ding{172} 是真命题， \ding{173} 是假命题}{\ding{172} 是假命题， \ding{173} 是真命题}{\ding{172}  \ding{173} 都是假命题}
    \begin{jieda}
        对于 \ding{172} ,数列 $\left\{  {b}_{n}\right\}$ 的前 $n$ 项和 ${T}_{n} = {\left( n - t\right) }^{2}(t$ 为正整数),当 $n \geq  2$ 时, ${b}_{n} = {T}_{n} - {T}_{n - 1} = {\left( n - t\right) }^{2} - {\left( n - 1 - t\right) }^{2} =  {2n} - {2t} - 1$ ,当 $n = 1$ 时, ${b}_{1} = {T}_{1} = {\left( 1 - t\right) }^{2}$ 不满足上式,所以 $\left| {b}_{n}\right|  = \left\{  {\begin{array}{l} {\left( 1 - t\right) }^{2},n = 1 \\  {2n} - {2t} - 1,n \geq  2 \end{array}\text{ ,当 }t = 1,n \geq  2}\right.$ 时, ${2n} - {2t} - 1 =  {2n} - 3 > 0$ ,所以数列 $\left\{  \left| {b}_{n}\right| \right\}$ 与原数列相同,所以 ${T}_{n} = (n -  1{)}^{2}$ ,所以当 $t = 1$ 时,数列 $\left\{  \left| {b}_{n}\right| \right\}$ 为完全平方数列,当 $t \geq  2$ 时, $\left| {b}_{1}\right|  + \left| {b}_{2}\right|  = {\left( t - 1\right) }^{2} + \left| {3 - {2t}}\right|  = {t}^{2} - {2t} + 1 + {2t} - 3 =  {t}^{2} - 2$ 不是 “完全平方数，所以当 $t \geq  2$ 时，数列 $\left\{  \left| {b}_{n}\right| \right\}$ 不是完全平方数列,综上所述: 只有当 $t = 1$ 时,数列 $\left\{  \left| {b}_{n}\right| \right\}$ 为 “完全平方数列”,故 \ding{172} 是真命题；
        
        对于 \ding{173} ，$a_n = 1$即可, 故  \ding{173}  是真命题. 
    \end{jieda}
    \item (2024$\cdot$黄浦区二模$\cdot$16)设数列 $\left\{  {a}_{n}\right\}$ 的前 $n$ 项和为 ${S}_{n}$ ，若对任意的正整数 $n$ ， ${S}_{n}$ 都是数列 $\left\{  {a}_{n}\right\}$ 中的项，则称数列 $\left\{  {a}_{n}\right\}$ 为“ $T$ 数列”. 对于命题: \ding{172}  存在“ $T$ 数列” $\left\{  {a}_{n}\right\}$ ，使得数列 $\left\{  {S}_{n}\right\}$ 为公比不为 1 的等比数列;  \ding{173}  对于任意的实数 ${a}_{1}$ ,都存在实数 $d$ ,使得以 ${a}_{1}$ 为首项、 $d$ 为公差的等差数列 $\left\{  {a}_{n}\right\}$ 为“ $T$ 数列”. 下列判断正确的是 \dotfill{(\kh{A})}
    \xx{\ding{172} 和 \ding{173} 均为真命题}{\ding{172} 和 \ding{173} 均为假命题}{\ding{172} 是真命题， \ding{173} 是假命题}{\ding{172} 是假命题， \ding{173} 是真命题} 
    \begin{jieda}
        \begin{dingautolist}{172}
            \item 设$S_n = a_1q^{n-1}$，$a_n = S_n - S_{n-1} = a_1q^{n-2}(q-1), n \geq 2$，故$a_1q^{n-2}(q-1) = a_1q^{m-1}$对任意的$m$，存在$n$使得方程有解. 故考虑$q = 2$，此时$n = m+1$，故$a_n = \left\{\begin{array}{l}
                a_1, n = 1\\
                a_1 \times 2^{n-2}, n \geq 2
            \end{array}\right.$满足要求.
            \item 当$a_1 \geq 0$时，$d  = 1$即可，当$a_1 < 0$，$d = -1$即可.
        \end{dingautolist}
    \end{jieda}
    \item (2025$\cdot$杨浦区一模$\cdot$16)设无穷数列 $\left\{  {a}_{n}\right\}$ 的前 $n$ 项和为 ${S}_{n}$ ，且对任意的正整数 $n$ ， ${a}_{n + 1} = \dfrac{{S}_{n}}{{a}_{n}}$ ，则 $\mathop{\sum }\limits_{{i = 1}}^{5}{a}_{2i} - \mathop{\sum }\limits_{{i = 1}}^{6}{a}_{{2i} - 1}$ 的值可能为 \dotfill{(\kh{A})}
    \xx{$-6$}{$0$}{$6$}{$12$}
    \begin{jieda}
        当 $n = 1$ 时, ${a}_{2} = \dfrac{{S}_{1}}{{a}_{1}} = 1$ . 当 $n \geq  2$ 时, ${S}_{n} = {a}_{n}{a}_{n + 1}$ , 所以 ${S}_{n - 1} = {a}_{n - 1}{a}_{n}$ ,两式相减得: ${a}_{n} = {a}_{n}\left( {{a}_{n + 1} - {a}_{n - 1}}\right)$ ,因为 ${a}_{n} \neq  0$ ,所以 ${a}_{n + 1} - {a}_{n - 1} = 1$ . 所以数列 $\left\{  {a}_{n}\right\}$ 的奇数项是以 ${a}_{1}$ 为首项,1 为公差的等差数列,且 ${a}_{1} \neq  0$ . 所以 $\mathop{\sum }\limits_{{i = 1}}^{6}{a}_{{2i} - 1} =  6{a}_{1} + \dfrac{6 \times  5}{2} \times  1 = 6{a}_{1} + {15}$ . 同理,数列 $\left\{  {a}_{n}\right\}$ 的偶数项是以 1 为首项，1 为公差的等差数列，所以 $\mathop{\sum }\limits_{{i = 1}}^{5}{a}_{2i} = 5 \times  1 + \dfrac{5 \times  4}{2} \times  1 = {15}$ . 所以 $\mathop{\sum }\limits_{{i = 1}}^{5}{a}_{2i} - \mathop{\sum }\limits_{{i = 1}}^{6}{a}_{{2i} - 1} = {15} - 6{a}_{1} - {15} =  - 6{a}_{1}$ . 若 $- 6{a}_{1} =  - 6 \Longrightarrow  {a}_{1} = 1$ ,则数列 $\left\{  {a}_{n}\right\}$ 各项均不为 0,数列 $\left\{  {a}_{n}\right\}$ 是无穷数列,故 ${A}$ 正确; 若 $- 6{a}_{1} = 0 \Longrightarrow  {a}_{1} = 0$ ,这与 ${a}_{1} \neq  0$ 矛盾,故 B 错误; 若 $- 6{a}_{1} = 6 \Longrightarrow  {a}_{1} =  - 1$ ,根据奇数项成公差为 1 的等差数列,则 ${a}_{3} = 0$ ,则 ${a}_{4}$ 无法求出,这与数列 $\left\{  {a}_{n}\right\}$ 是无穷数列矛盾,故 ${C}$ 错误; 若 $- 6{a}_{1} = {12} \Longrightarrow  {a}_{1} =  - 2$ ,根据奇数项成公差为 1 的等差数列,则 ${a}_{5} = 0$ ,则 ${a}_{6}$ 无法求出,这与数列 $\left\{  {a}_{n}\right\}$ 是无穷数列矛盾,故 ${D}$ 错误. 故选: A.
    \end{jieda}
\end{enumerate}
\subsection*{$S_n, S_{2n}-S_n, S_{3n}-S_{2n}, ...$}
\begin{enumerate}
    \item 记 ${S}_{n}$ 为等差数列 $\left\{  {a}_{n}\right\}$ 的前 $n$ 项和,若 $\dfrac{{S}_{4}}{{S}_{8}} = \dfrac{1}{3}$ ,则 $\dfrac{{S}_{8}}{{S}_{16}} =$\blkda{$\dfrac{3}{10}$}
\end{enumerate}
\subsection*{通项公式与数列的单调性}
\begin{enumerate}
    \item (2025$\cdot$嘉定区一模$\cdot$8)已知数列 $\left\{  {a}_{n}\right\}$ 的通项公式为 ${a}_{n} =  - n + c$ ，其中 $c$ 为常数，设数列 $\left\{  {a}_{n}\right\}$ 的前 $n$ 项和为 ${S}_{n}$ ，若 ${S}_{6} > {S}_{5}$ 且 ${S}_{6} > {S}_{7}$ ，则 $c$ 的取值范围为\blkda[2]{$(6, 7)$}.
    \begin{jieda}
        ${S}_{6} > {S}_{5} \Longleftrightarrow a_6 > 0$，${S}_{6} > {S}_{7} \Longleftrightarrow a_7 < 0$，$\left\{\begin{array}{l}
            c-6 > 0 \\
            c-7 < 0
        \end{array}\right.$，故$c \in (6, 7)$.
    \end{jieda}

    \item (2025$\cdot$上海秋季卷$\cdot$16)已知数列 $\left\{  {a}_{n}\right\}$ 、 $\left\{  {b}_{n}\right\}$ 、 $\left\{  {c}_{n}\right\}$ 的通项公式分别为 ${a}_{n} = {10n} - 9,{b}_{n} = {2}^{n},{c}_{n} = \lambda {a}_{n} +  \left( {1 - \lambda }\right) {b}_{n}$. 若对任意的 $\lambda  \in  \left\lbrack  {0,1}\right\rbrack$ ， ${a}_{n}$ 、 ${b}_{n}$ 、 ${c}_{n}$ 的值均能构成三角形，则满足条件的正整数 $n$ 有\dotfill{(\kh{B})}
    \xx{4 个}{3 个}{1 个}{无数个}
    \begin{jieda}
        由题意可知$min\{a_n, b_n\} \times 2 > max\{a_n, b_n\}$，易知$\{b_n\}$增长快于$\{a_n\}$，故枚举前几项即可.
    \end{jieda}
\end{enumerate}
\subsection*{递推公式与数列的单调性}
\begin{enumerate}
    \item (2025$\cdot$闵行区一模$\cdot$16)已知数列 $\left\{  {a}_{n}\right\}$ 满足 ${a}_{n + 1} = \left| {{a}_{n} + 1}\right|  + \lambda \left| {{a}_{n} - 1}\right|$ ，其中 $\lambda$ 为常数. 对于下述两个命题:

    \ding{172}  对于任意的 $\lambda  > 0$ ，任意的 ${a}_{1} \in  \mathbf{R}$ ，都有 $\left\{  {a}_{n}\right\}$ 是严格增数列；

    \ding{173}  对于任意的 $\lambda  < 0$ ，存在 ${a}_{1} \in  \mathbf{R}$ ，使得 $\left\{  {a}_{n}\right\}$ 是严格减数列.

    以下说法正确的为 \dotfill{(\kh{A})}

    \xx{\ding{172} 真命题， \ding{173} 假命题}{\ding{172} 假命题， \ding{173} 真命题}{\ding{172}  \ding{173} 均为真命题}{\ding{172}  \ding{173} 均为假命题} 
    \begin{jieda}
        \begin{dingautolist}{172}
            \item 当$\lambda > 0$时，易知$a_n > 0, n > 1$，而$y = |x+1|+\lambda|x-1|$没有不动点，根据图象可知严增.
            \item 当$\lambda = -1$时，根据$y = |x+1|-|x-1|$的图像结合不动点分类讨论可知命题为假.
        \end{dingautolist}
    \end{jieda}

    \item (2025$\cdot$嘉定区一模$\cdot$16)已知数列 $\left\{  {a}_{n}\right\}$ 满足 ${a}_{n + 1} = r{a}_{n}\left( {1 - {a}_{n}}\right) \left( {n = 1,2,3,\cdots }\right)$ ， ${a}_{1} \in  \left( {0,1}\right)$ ，给出以下四个结论:

    \ding{172}  当 $r = 2$ 时，只存在有限个 ${a}_{1}$ ，使得对任意正整数 $n$ ，都有 ${a}_{n + 1} > {a}_{n}$ ；

    \ding{173}  当 $r = 2$ 时，存在 ${a}_{1}$ 和正整数 $P$ ，当 $n > P$ 时， ${a}_{n + 1} - {a}_{n} < \dfrac{1}{2025}$ ；

    \ding{174}  当 $r = 3$ 时，存在 ${a}_{1}$ 和正整数 $P$ ，当 $n > P$ 时， ${a}_{n + 1} = {a}_{n}$ ；

    \ding{175}  当 $r =  - 3$ 时,不存在 ${a}_{1}$ ,使得对任意正整数 $n$ ,且 $n \geq  3$ ,都有 ${a}_{n} > 0$. 
    
    其中正确结论是 \dotfill{(\kh{B})}

    \xx{\ding{172}  \ding{173}}{\ding{173}  \ding{174}}{\ding{174}  \ding{175}}{\ding{173}  \ding{175}}
    \begin{jieda}
        数形结合，结合不动点判断即可.
    \end{jieda}
\end{enumerate}
\subsection*{递推求通项}
\begin{enumerate}
    \item (2024$\cdot$普陀区一模$\cdot$9)若数列 $\left\{  {a}_{n}\right\}$ 满足 ${a}_{1} = {12},{a}_{n + 1} = {a}_{n} + {2n}$ ( $n \geq  1$ ， $n \in  \mathbf{N}$)，则 $\dfrac{{a}_{n}}{n}$ 的最小值是\blkda[2]{6}.
    \begin{jieda}
        对${a}_{n + 1} = {a}_{n} + {2n}$累加可得$a_n = n^2 - n + 12$，故$\dfrac{a_n}{n} = n + \dfrac{12}{n} - 1$，在$n = 3$或$n = 4$时取得最小值$6$.
    \end{jieda}
\end{enumerate}